\uuid{NQkj}
\exo7id{3598}
\titre{exo7 3598}
\auteur{quercia}
\organisation{exo7}
\datecreate{2010-03-10}
\isIndication{false}
\isCorrection{true}
\chapitre{Réduction d'endomorphisme, polynôme annulateur}
\sousChapitre{Autre}
\cpge{mpsi-mp2i.calcul-matriciel-systemes.a, pcsi.calcul-matriciel-systemes-lineaires.a, ptsi.calcul-matriciel-systemes-lineaires.a, mpsi-mp2i.calcul-matriciel-systemes.e, pcsi.determinants.c, ptsi.determinants.a, mpsi-mp2i.changements-bases.a, pcsi.matrices-applications-lineaires.c, ptsi.matrices-applications-lineaires.c, bcpst1.matrices, mp-mpi.reduction-endomorphismes.d, pc.reduction.c, psi.reduction-endomorphismes.c, pt.reduction-endomorphismes.c, mp-mpi.reduction-endomorphismes.i, pc.reduction.d, psi.reduction-endomorphismes.d, bcpst2.valeurs-propres-vecteurs-propres.b}
\module{Algèbre}
\niveau{L2}
\difficulte{}

\contenu{
\question{
Soit $M \in \mathcal{M}_n(\C)$ diagonalisable.
    Soit $A = \begin{pmatrix}M&M\cr M&M\cr\end{pmatrix} \in \mathcal{M}_{2n}(\C)$.
    La matrice $A$ est-elle diagonalisable ?
}
\reponse{
S'inspirer du cas $n=1$. Soit $P = \begin{pmatrix}I&I\cr I&-I\cr\end{pmatrix}$ :
	     $P^{-1}AP = \begin{pmatrix}2M&0\cr 0&0\cr\end{pmatrix}$ est diagonalisable,
	     donc $A$ aussi.
}
}
